
We have proposed a simple, effective, and efficient approach for out-of-scope intent detection by overcoming the limitation of previous methods via matching train-test conditions. Particularly, at the training stage, we construct self-supervised and open-domain outliers to improve model generalization and simulate real outliers in the test stage. Extensive experiments on four dialogue datasets show that our approach significantly outperforms state-of-the-art methods. In the future, we plan to investigate the theoretical underpinnings of our approach and apply it to more applications.

%for out-of-scope intent detection in different scenarios such as cross-domain and few-shot learning scenarios.
%, and achieves large performance gains especially when the amount of training data is relatively small. I

%forming a consistent $K+1$ classification task for training and test. In particular, we construct \emph{hard} self-supervised outliers and \emph{easy} open-domain outliers at the training stage without utilizing any of information of outlliers in the test set. Extensive experiments demonstrate the effectiveness of our method. Our method results in favorable performance on unknown intent classification and provides especially significant improvement in data-scarce scenarios. In the future, we will study the effect of our propose method in diverse unknown intent classification tasks, such as cross-domain and few-shot scenarios.